\section{Introdução}

\IEEEPARstart{A}{} série de Fourier representa um sinal periódico como a soma de uma
componente contínua e de senos e cossenos em múltiplos da frequência
fundamental. Essa decomposição permite reconstruir formas de onda,
comparar seu conteúdo espectral e quantificar a contribuição energética
de cada harmônico.

Para comprovar esse comportamento, serão utilizados diversos sinais e calculados
suas representações em série de Fourier e sua energia. O desenvolvimento
preserva as amplitudes e os intervalos como variáveis e apresenta as
integrais e suas simplificações. Em seguida, trata os deslocamentos e a
mudança do período, calcula as energias por Parseval e compara os
resultados com a FFT. 

